% SHARD-ID: AQM-01-PROGRAMME-MAP
% SHARD-TITLE: Programme Map
% SHARD-SUMMARY: Organizes the intended research axes without asserting an established bridge between them.
% SHARD-SUMMARY: Lists the first convention and source decisions needed before technical derivations begin.
% SHARD-KEYWORDS: programme, Weil conjectures, zeta functions, arithmetic QM, SUSY QM, toric code, Levin-Wen

\section{Programme Map}

This shard is a map of questions, not a list of results.

\subsection{Arithmetic And Zeta-Function Side}

\begin{itemize}
  \item Which zeta functions are in scope: local zeta functions over finite
  fields, Hasse-Weil zeta functions, motivic zeta functions, or other
  variants?
  \item Which Frobenius convention and cohomology theory will be used for trace
  formulas?
  \item Which finite examples are small enough for exact Sage or GAP checks?
\end{itemize}

\subsection{Quantum-Mechanical Side}

\begin{itemize}
  \item What is meant locally by arithmetic quantum mechanics: a literal
  Hilbert-space model, a spectral dictionary, or a finite-dimensional
  approximation?
  \item Which supersymmetric-QM convention fixes the grading, supercharge,
  adjoint, Hamiltonian, and Witten-index normalization?
  \item Which invariants can be checked before proposing an analogy:
  nilpotence, positivity, trace identities, index stability, or spectral
  reconstruction?
\end{itemize}

\subsection{Topological-Code Side}

\begin{itemize}
  \item Which lattice model is being used: toric code, quantum double, or
  Levin-Wen string-net?
  \item Which boundary conditions, lattice orientation, and star/plaquette
  conventions are fixed?
  \item Which category data are required: fusion rules, F-symbols, R-symbols,
  pivotal or spherical structure, and Drinfeld-center data?
\end{itemize}

\subsection{First Decisions}

The first technical session should not start with formulas. It should register
sources and fix the minimum conventions needed to compare one finite
arithmetic zeta example with one finite quantum or topological model invariant.
